\subsection{Motivación: no linealidad con semántica agronómica}

El Checkpoint~03A introduce una capa aditiva de 350 variables derivadas sin utilizar
\artifact{RENDIMIENTO_T_HA}. El objetivo no es reemplazar Feature Table v1, sino codificar
hipótesis no lineales interpretables antes de abrir una búsqueda de modelos. Con aproximadamente
1,400 predictores base, una expansión polinómica completa de segundo orden generaría del orden
de $10^6$ interacciones; esa estrategia es inadecuada para $n=138$. En su lugar, se construyen
familias pequeñas cuya razón agronómica puede enunciarse antes de observar $y$.

El build final contiene 197 filas y 350 features: 123 clean y 227 competition. No hay missing
ni infinitos. Tres variables de calor resultaron constantes y se eliminaron automáticamente.
Las familias retenidas son:
\[
\begin{array}{lr}
\text{fenología} & 162\\
\text{anomalías fenológicas} & 41\\
\text{agua/productividad} & 40\\
\text{térmicas} & 36\\
\text{cross-domain} & 19\\
\text{perfil de suelo} & 18\\
\text{timing de agua} & 10\\
\text{sensor agreement} & 9\\
\text{interacciones de suelo} & 8\\
\text{bases no lineales} & 7
\end{array}
\]

\subsection{Fenología y forma estacional}

Para una trayectoria mensual $x_t$, $t\in\{4,\ldots,10\}$, se calculan media estacional,
amplitud, AUC, mes de máximo, medias temprana/tardía, cambio junio--agosto y pendientes de
green-up/senescence. Por ejemplo,
\begin{align}
\bar x &= \frac{1}{T}\sum_{t=1}^{T}x_t,\\
A(x) &= \max_t x_t-\min_t x_t,\\
AUC(x) &\approx
\sum_{t=1}^{T-1}\frac{x_t+x_{t+1}}{2}
(m_{t+1}-m_t),\\
t_{\max} &= \arg\max_t x_t.
\end{align}

Las pendientes son coeficientes de regresiones lineales locales sobre bloques de meses:
\[
x_t=\alpha+\beta t+\varepsilon_t.
\]
No se afirma que la fenología real sea lineal; $\beta$ funciona como descriptor compacto de
dirección de cambio. La literatura de cebada y estimación de rendimiento por teledetección
motiva conservar timing y forma en lugar de depender sólo de una media
\citep{Chanev2025,Sharifi2020,Mirosavljevic2018}.

\subsection{Anomalías 2025 frente a historia}

Una parcela puede tener productividad estructuralmente alta o baja y, al mismo tiempo, mostrar
una condición 2025 inusual respecto a su propia historia. La capa separa ambas dimensiones
mediante diferencias y razones:
\begin{align}
\Delta\bar x_i
&=
\bar x_{i,2025}-\bar x_{i,\mathrm{hist}},\\
R_{x,i}
&=
\frac{\bar x_{i,2025}}
{\bar x_{i,\mathrm{hist}}+\epsilon},\\
\Delta AUC_i
&=
AUC_{i,2025}-AUC_{i,\mathrm{hist}}.
\end{align}

También se incluyen delta de agosto, delta de late-season y desplazamiento del mes de pico.
Para NDWI se evita deliberadamente una razón de medias, porque el denominador puede cruzar
cero al tratarse de un índice signed. Esta es una decisión numérica y semántica: una feature
de gran magnitud causada por un denominador cercano a cero no debe confundirse con una señal
agronómica extrema.

\subsection{Concordancia inter-sensor}

Sentinel-2 y Planet 2025 se comparan para NDVI, EVI y LAI mediante diferencia de media,
razón de media y RMSE mensual:
\begin{equation}
\RMSE_{\mathrm{sensor},i}
=
\sqrt{
\frac{1}{T}\sum_t
\left(x^{S2}_{i,t}-x^{Planet}_{i,t}\right)^2
}.
\label{eq:sensor-rmse}
\end{equation}
La feature no presupone equivalencia de nivel entre productos. Su interpretación es
diagnóstica: si dos sensores describen de manera muy distinta la misma parcela, esa
discordancia puede contener información de cobertura, procesamiento, nubosidad o dinámica
vegetal.

\subsection{Desarrollo térmico}

A partir de temperatura máxima y mínima se construye
\[
T_{\mathrm{mean},m}=\frac{T_{\max,m}+T_{\min,m}}{2},
\qquad
DTR_m=T_{\max,m}-T_{\min,m}.
\]
Se agregan proxies mensuales de growing degree days:
\begin{equation}
GDD_{\mathrm{proxy}}(T_b)
=
\sum_m d_m\max(T_{\mathrm{mean},m}-T_b,0),
\label{eq:gdd-proxy}
\end{equation}
para $T_b=0^\circ$C y $5^\circ$C. Se usa explícitamente la palabra proxy porque los inputs
climáticos son mensuales, no diarios. La acumulación térmica es una construcción habitual en
fenología de cereales \citep{McMaster1997,McMaster2005,Hajkova2019}, pero el valor exacto aquí
no debe confundirse con un GDD agronómico diario de alta resolución.

Se propuso también
\[
H_{25}
=
\sum_m d_m\max(T_{\mathrm{mean},m}-25,0).
\]
En la build canónica resultó constante tanto histórica como contemporáneamente y fue eliminado,
junto con la interacción que lo utilizaba. Este comportamiento es útil: el pipeline no fuerza
a conservar una hipótesis sólo porque conceptualmente sea razonable.

\subsection{Timing de precipitación}

La precipitación se separa en ventanas:
\begin{align}
P_{\mathrm{early}} &= P_{\mathrm{Apr}}+P_{\mathrm{May}}+P_{\mathrm{Jun}},\\
P_{\mathrm{mid}} &= P_{\mathrm{Jun}}+P_{\mathrm{Jul}}+P_{\mathrm{Aug}},\\
P_{\mathrm{late}} &= P_{\mathrm{Aug}}+P_{\mathrm{Sep}}+P_{\mathrm{Oct}}.
\end{align}
También se calculan coeficiente de variación y fracción tardía. La motivación es que la
respuesta del cultivo a déficit hídrico depende de la etapa de desarrollo
\citep{Bello2022,Mogensen1980,Zhang2015}. El total estacional no contiene por sí solo
información sobre cuándo ocurrió el agua.

\subsection{Productividad hídrica y WaPOR}

AETI y NPP se combinan en:
\begin{align}
WP_{\mathrm{NPP}}
&=
\frac{NPP}{AETI+\epsilon},\\
R_{ET/P}
&=
\frac{AETI}{P+\epsilon},\\
WB
&=
P-AETI.
\end{align}
La primera es un proxy de productividad de biomasa por agua observada desde teledetección, no
una definición de water-use efficiency de grano. De forma similar, $AETI/P$ no es un índice de
sequía: el cociente puede estar influido por almacenamiento de agua en suelo, riego u otras
fuentes. La nomenclatura deliberadamente evita atribuir causalidad fisiológica a una
transformación predictiva \citep{Dhonthi2024}.

Se añaden interacciones como
\[
NPP\times NDVI,\qquad LAI\times AETI,
\]
para permitir a modelos lineales expresar que la biomasa y el uso de agua pueden tener
efectos dependientes del estado del dosel.

\subsection{Perfil edáfico y bases no lineales}

Para cada propiedad $p$ de SoilGrids se construyen contraste vertical y razón:
\begin{align}
G_p &= p_{0-30}-p_{30-60},\\
R_p &= \frac{p_{0-30}}{p_{30-60}+\epsilon}.
\end{align}
Adicionalmente se crean interacciones compactas: SOC$\times$CEC,
SOC$\times$N, clay$\times$SOC, clay$\times$CEC, bulk-density$\times$clay,
clay/sand, clay+silt y $pH^2$. La forma cuadrática de pH no fija un óptimo; simplemente
permite que un modelo lineal represente curvatura junto con el término lineal. La respuesta de
cebada a química y cationes de suelo puede ser no lineal \citep{Holland2021}.

Variables positivas y sesgadas reciben una base
\[
z=\log(1+x),
\]
principalmente para estabilizar escalas en modelos lineales.

\subsection{Interacciones cross-domain}

Se crean hipótesis entre dominios, entre ellas:
\[
Y_{\mathrm{SIAP,hist}}\times\Delta NDVI_{2025},
\qquad
Y_{\mathrm{SIAP,hist}}\times NPP_{2025},
\]
así como slope$\times$precipitación, elevación$\times$temperatura,
clay$\times$precipitación, SOC$\times$precipitación, N de suelo$\times$NDVI,
NDVI$\times$NDWI y EVI$\times$LAI. El objetivo es proporcionar bases físicamente
interpretables sin enumerar todas las interacciones posibles.

\subsection{Clases de evidencia}

El manifest asigna a cada feature una clase:
\begin{description}
\item[literature-backed] la relación general está motivada por literatura;
\item[mechanistic proxy] la transformación tiene interpretación física/agronómica clara,
pero la fórmula exacta es un proxy del proyecto;
\item[experimental] base matemática target-free que debe demostrar valor predictivo.
\end{description}

Esta clasificación impide escribir retrospectivamente que una interacción que funcionó en CV
es una ley agronómica. La validez predictiva se decide en Checkpoint~03C; la validez
mecanicista no se infiere del score.
